\section{Технические документации}

\subsection{Организация репозиториев}

Разработка системы ведётся в двух отдельных репозиториях: один для серверной (бэкенд) части, другой — для клиентской (фронтенд) части приложения.

Бэкенд-репозиторий содержит серверную часть приложения: реализацию игровой логики и правил, серверный API, механизмы хранения и обработки состояния игровых сессий, управление пользователями и ролями, а также синхронизацию данных между клиентами в реальном времени.

Фронтенд-репозиторий содержит клиентскую часть приложения: пользовательский интерфейс для участников и ведущего, визуализацию состояния игры и метрик, а также логику взаимодействия с сервером через API.

Разделение кодовой базы на два репозитория обусловлено следующими соображениями:

\begin{itemize}
    \item \textbf{Независимость разработки}: команды фронтенда и бэкенда могут работать параллельно без риска конфликтов при слиянии изменений.
    \item \textbf{Разграничение ответственности}: каждый репозиторий имеет собственный цикл сборки, тестирования и деплоя, что упрощает настройку CI/CD-пайплайнов.
    \item \textbf{Технологическая независимость}: стек технологий и зависимости фронтенда и бэкенда не пересекаются, что снижает сложность управления проектом.
    \item \textbf{Контроль доступа}: при необходимости можно гибко управлять правами доступа участников команды к каждой части системы.
\end{itemize}

Взаимодействие между двумя частями системы осуществляется через сетевой API, документация к которому является общим артефактом и поддерживается совместно обеими командами.
